\section{Limitations and Threats to Validity}
\label{sec:limitations}

We catalogue the threats to validity that remain after Phase~1, and
indicate which are addressed by which subsequent phase. Phase~2 is
the binding open work.

\subsection{Sample Size}

Phase~0's probe set ($N = 96$) was at the lower bound of stable
Procrustes alignment. Phase~1 raised $N$ to $600$ and the per-domain
stratification used $n = 200$ each. Both regimes satisfy the
$N \geq 5k$ rule of thumb at the operative dimensions and produce
$z$-scores far above the threshold at which sample-size-driven null
inflation could plausibly change the verdict.

\subsection{Probe Set Domain Coverage}

Phase~0 tested a single domain (political-economy / categorial-architecture).
Phase~1 extended coverage to three distinct compositional-reasoning
domains (mathematics, logic, political-economy), with per-domain
stratification confirming the alignment signal in each domain
independently. The COSI hypothesis as stated in
Proposition~\ref{prop:cosi} is domain-general; the present empirical
support is domain-plural. Genuine cross-domain generality (e.g.,
natural images, audio, multimodal embeddings) remains untested.

\subsection{Probe Length}

Probes remain short (typically $12$--$30$ tokens) in both phases.
Last-token pooling on short sequences may bias the activation toward
prompt-template syntax rather than proposition content. The pooling
sweep (mean, max) was deferred from Phase~1 due to compute scheduling
and is committed for Phase~1b. Longer-probe replication ($30$--$200$
tokens) is deferred to the same follow-up.

\subsection{Shared Training Distribution}

This is the most important confound and the one Phase~1 does \emph{not}
address. Both Phi-4-Reasoning-Plus and Qwen3-Next-80B-A3B-Thinking are
trained on post-2024 web corpora that overlap substantially. The
aligned subspace we identify could plausibly be a residue of shared
training-distribution statistics rather than a substrate-independent
compositional structure. The cleanest test would replace one of the two
models with an architecturally similar but training-distribution-distant
model, holding the other fixed and observing whether the residual
remains low.

A near-perfect such control would be a model trained on a substantially
disjoint corpus (e.g., a non-web-pretrained scientific-corpus model, or a
multilingual model whose training distribution is dominated by a different
language). Such a model is not in our local zoo. Phase~2 of the program
is structured around this control. We flag this as the binding remaining
work.

\subsection{Architectural Distance Within the Pair}

Phase~0 used two dense pure-attention transformers (Phi-4 and Qwen3-4B).
Phase~1 paired the dense pure-attention Phi-4 with the hybrid
attention/state-space-model Qwen3-Next-80B, materially increasing the
architectural distance and confirming that the COSI signal survives the
attention/SSM class boundary. Architectures further afield (pure SSM,
recurrent, retrieval-augmented, multimodal) remain untested; Phase~4
extends the test across this broader architectural landscape.

\subsection{Linearity of the Alignment}

Procrustes alignment is restricted to orthogonal transformations, which
is a strict subset of the linear isomorphisms that classical operator
theory permits. The COSI hypothesis as stated in
Proposition~\ref{prop:cosi} requires only \emph{some} approximate
isomorphism between the subspaces; the orthogonal restriction may
underestimate the alignment if the true correspondence involves
anisotropic scaling. We compute CKA in Phase~1 as a complementary metric
that is invariant to isotropic scaling but not to anisotropic scaling, and
will report any discrepancy between the two verdicts.

We additionally note that the COSI hypothesis is in principle compatible
with non-linear isomorphisms that no linear method recovers. Establishing
such isomorphisms empirically requires either a generative-model-based
alignment method or a stronger theoretical bridge; we treat this as
beyond the scope of the present paper.

\subsection{Correspondence vs.\ Causation}

A positive COSI result is correspondence evidence: the two models'
representations are geometrically aligned on a shared probe set. It is
not causation evidence: it does not establish that the aligned subspace
is what the models \emph{use} to compute the response, as opposed to a
statistical shadow of something they share. Causal evidence requires
intervention: ablate a direction in $V_A$, predict the corresponding
ablation effect in $R(V_A) \subseteq V_B$, and run both. We treat
intervention experiments as Phase~3, requiring causal-tracing
infrastructure beyond the scope of the present calibration.

\subsection{Pooling and Layer Sensitivity}

The Phase~0 result uses last-token pooling at three fixed layer fractions.
Both choices may interact with the result in ways not yet characterized.
Phase~1 includes the full pooling and layer-fraction sensitivity sweeps
declared in the design note.
